\documentclass[12pt,a4paper]{article}
\usepackage[utf8]{inputenc}
\usepackage[russian]{babel}
\usepackage[T2A]{fontenc}
\usepackage{geometry}
\geometry{left=2.5cm,right=1.5cm,top=2cm,bottom=2cm}
\usepackage{amsmath,amssymb}
\usepackage{graphicx}
\usepackage{tikz}
\usetikzlibrary{positioning,calc}
\usepackage{listings}
\usepackage{xcolor}
\usepackage{hyperref}

% Настройка листингов кода
\lstset{
  language=[LaTeX]TeX,
  basicstyle=\small\ttfamily,
  keywordstyle=\color{blue},
  commentstyle=\color{green!60!black},
  stringstyle=\color{red},
  breaklines=true,
  frame=single,
  numbers=left,
  numberstyle=\tiny\color{gray},
  backgroundcolor=\color{gray!5}
}

\title{\textbf{Отчет по лабораторной работе №8} \\ 
       \large Построение графических объектов в TikZ}
\author{Выполнил: [Викторов Егор] \\ Группа: [НПМмд02-24]}
\date{\today}

\begin{document}

\maketitle
\tableofcontents
\newpage

\section{Цель работы}

Целью данной работы является изучение возможностей пакета TikZ для создания различных графических объектов в системе LaTeX, включая:
\begin{itemize}
  \item построение графов с узлами различных стилей;
  \item визуализацию математических функций;
  \item реализацию рекурсивных алгоритмов для построения фракталов.
\end{itemize}

\section{Теоретические сведения}

\subsection{Пакет TikZ}

TikZ (TikZ ist kein Zeichenprogramm) — мощный инструмент для создания векторной графики непосредственно в документах LaTeX. Основные преимущества:
\begin{itemize}
  \item Интеграция с математическими формулами
  \item Точность и масштабируемость
  \item Программируемость и возможность создания макросов
  \item Поддержка сложных геометрических построений
\end{itemize}

\subsection{Фрактал «Ковер Серпинского»}

Ковер Серпинского — фрактальная фигура, получаемая рекурсивным удалением центральных квадратов. На каждом этапе квадрат делится на 9 равных частей (сетка 3×3), и центральная часть удаляется. Процесс повторяется для оставшихся 8 квадратов.

Фрактальная размерность ковра Серпинского:
\[
D = \frac{\ln 8}{\ln 3} \approx 1.8928
\]

\section{Выполнение работы}

\subsection{Задание 1: Построение графа}

\subsubsection{Описание задачи}

Требовалось построить граф с 6 узлами, расположенными по окружности. Узлы имеют два стиля:
\begin{itemize}
  \item \textbf{green} — зеленые заполненные круги (узлы A, C, E)
  \item \textbf{ring} — узлы с двойной обводкой (узлы B, D, F)
\end{itemize}

Граф содержит различные типы ребер:
\begin{itemize}
  \item Пунктирные синие линии с весами (B--F: 6, B--D: 2, F--D: 4)
  \item Голубая кривая с закругленными углами (B--C--D)
  \item Красные гладкие кривые (A--D, A--E, A--B)
\end{itemize}

\subsubsection{Реализация}

\begin{lstlisting}
\documentclass[border=1cm]{standalone}
\usepackage{tikz}
\usetikzlibrary{positioning,calc}
\begin{document}
\begin{tikzpicture}[>=stealth, every node/.style={font=\small}]
  \tikzset{
    ring/.style={circle,draw=black,thin,double=black,
                 double distance=1.5pt,minimum size=9mm,
                 inner sep=0pt,fill=white},
    green/.style={circle,draw=black,very thick,
                  minimum size=9mm,inner sep=0pt,fill=green!60},
  }

  % Узлы в полярных координатах
  \node[green] (A) at (90:3cm) {A};
  \node[ring]  (B) at (150:3cm) {B};
  \node[green] (C) at (210:3cm) {C};
  \node[ring]  (D) at (270:3cm) {D};
  \node[green] (E) at (330:3cm) {E};
  \node[ring]  (F) at (30:3cm)  {F};

  % Пунктирные линии с метками
  \draw[dotted,blue,line width=1.5pt] (B) -- 
    node[midway,above,fill=white] {6} (F);
  \draw[dotted,blue,very thick] (B) -- 
    node[pos=0.5,left,fill=white] {2} (D);
  \draw[dotted,blue,very thick] (F) -- 
    node[pos=0.5,right,fill=white] {4} (D);

  % Голубая кривая
  \draw[cyan!80!black,very thick,rounded corners=6pt]
    (B) to[out=200,in=160] (C) to[out=-20,in=180] (D);

  % Красные кривые
  \draw[red!80!black,very thick,smooth]
    (A) to[out=-60,in=90]
    node[pos=0.5,right,red,fill=white] {$\sqrt{2}$} (D);
  \draw[red!80!black,very thick,smooth]
    (A) to[out=-30,in=90] (E);
  \draw[red!80!black,very thick,smooth]
    (A) to[out=180,in=90] (B);
\end{tikzpicture}
\end{document}
\end{lstlisting}

\subsubsection{Результат}

\begin{center}
\begin{tikzpicture}[>=stealth, every node/.style={font=\small},scale=0.7]
  \tikzset{
    ring/.style={circle,draw=black,thin,double=black,double distance=1.5pt,minimum size=9mm,inner sep=0pt,fill=white},
    green/.style={circle,draw=black,very thick,minimum size=9mm,inner sep=0pt,fill=green!60},
  }
  \node[green] (A) at (90:3cm) {A};
  \node[ring]  (B) at (150:3cm) {B};
  \node[green] (C) at (210:3cm) {C};
  \node[ring]  (D) at (270:3cm) {D};
  \node[green] (E) at (330:3cm) {E};
  \node[ring]  (F) at (30:3cm)  {F};
  \draw[dotted,blue,line width=1.5pt] (B) -- node[midway,above,fill=white] {6} (F);
  \draw[dotted,blue,very thick] (B) -- node[pos=0.5,left,fill=white] {2} (D);
  \draw[dotted,blue,very thick] (F) -- node[pos=0.5,right,fill=white] {4} (D);
  \draw[cyan!80!black,very thick,rounded corners=6pt]
    (B) to[out=200,in=160] (C) to[out=-20,in=180] (D);
  \draw[red!80!black,very thick,smooth]
    (A) to[out=-60,in=90] node[pos=0.5,right,red,fill=white] {$\sqrt{2}$} (D);
  \draw[red!80!black,very thick,smooth]
    (A) to[out=-30,in=90] (E);
  \draw[red!80!black,very thick,smooth]
    (A) to[out=180,in=90] (B);
\end{tikzpicture}
\end{center}

\subsection{Задание 2: Построение графиков функций}

\subsubsection{Описание задачи}

Необходимо построить графики двух взаимно обратных функций:
\begin{itemize}
  \item Экспоненциальная функция: $y = e^x$
  \item Натуральный логарифм: $y = \ln(x)$
\end{itemize}

График должен включать координатные оси, вспомогательную линию $y=1$, метки и подписи.

\subsubsection{Реализация}

\begin{lstlisting}
\documentclass[border=1cm]{standalone}
\usepackage{tikz}
\begin{document}
\begin{tikzpicture}[x=1.2cm,y=1.2cm]
  % Координатные оси
  \draw[->,gray!60] (0,-2.6) -- (0,3.6) node[above] {$y$};
  \draw[->,gray!60] (-1.6,0) -- (3.2,0) node[right] {$x$};

  % Вспомогательная линия y=1
  \draw[gray!25] (-1.3,1) -- (3.0,1);
  \draw[gray!60] (-0.05,1) -- (0.05,1);
  \node[gray!50,left] at (-1.1,1) {$y=1$};

  % Открытый круг в начале координат
  \draw[fill=white,draw=gray!60] (0,0) circle(0.06);

  % Метка x=1
  \draw[gray!60] (1,-0.05) -- (1,0.05);
  \node[gray!40,below=3pt] at (1,0) {$x=1$};

  % График y = e^x
  \draw[domain=-1.2:2.2,smooth,variable=\x,
        blue!80!black,very thick]
    plot ({\x},{exp(\x)});
  \node[blue!80!black, right] at (2.05,3.05) {$y=e^{x}$};

  % График y = ln(x)
  \draw[domain=0.11:2.8,smooth,variable=\x,
        black!80,very thick]
    plot ({\x},{ln(\x)});
  \node[black!80,right] at (2.3,0.85) {$y=\ln(x)$};
\end{tikzpicture}
\end{document}
\end{lstlisting}

\subsubsection{Результат}

\begin{center}
\begin{tikzpicture}[x=1.2cm,y=1.2cm,scale=0.8]
  \draw[->,gray!60] (0,-2.6) -- (0,3.6) node[above] {$y$};
  \draw[->,gray!60] (-1.6,0) -- (3.2,0) node[right] {$x$};
  \draw[gray!25] (-1.3,1) -- (3.0,1);
  \draw[gray!60] (-0.05,1) -- (0.05,1);
  \node[gray!50,left] at (-1.1,1) {$y=1$};
  \draw[fill=white,draw=gray!60] (0,0) circle(0.06);
  \draw[gray!60] (1,-0.05) -- (1,0.05);
  \node[gray!40,below=3pt] at (1,0) {$x=1$};
  \draw[domain=-1.2:2.2,smooth,variable=\x,blue!80!black,very thick]
    plot ({\x},{exp(\x)});
  \node[blue!80!black, right] at (2.05,3.05) {$y=e^{x}$};
  \draw[domain=0.11:2.8,smooth,variable=\x,black!80,very thick]
    plot ({\x},{ln(\x)});
  \node[black!80,right] at (2.3,0.85) {$y=\ln(x)$};
\end{tikzpicture}
\end{center}

\subsection{Задание 3: Ковер Серпинского}

\subsubsection{Описание задачи}

Реализовать рекурсивный алгоритм построения ковра Серпинского и визуализировать первые 5 этапов (от 0 до 4).

\subsubsection{Алгоритм}

Рекурсивная функция \texttt{\textbackslash carpet\{x\}\{y\}\{size\}\{level\}}:
\begin{enumerate}
  \item \textbf{Базовый случай} (level = 0): рисуется заполненный квадрат
  \item \textbf{Рекурсивный случай}:
  \begin{itemize}
    \item Квадрат делится на 9 частей (сетка 3×3)
    \item Центральная часть пропускается
    \item Для остальных 8 частей рекурсивно вызывается функция с уровнем level-1
  \end{itemize}
\end{enumerate}

\subsubsection{Реализация}

\begin{lstlisting}
\documentclass[border=2mm]{standalone}
\usepackage{tikz}
\begin{document}

\newcommand{\carpet}[4]{%
  \pgfmathtruncatemacro{\level}{#4}%
  \ifnum\level=0
    \fill[black] (#1,#2) rectangle +(#3,#3);
  \else
    \pgfmathsetmacro{\newsize}{#3/3}%
    \pgfmathtruncatemacro{\nextlevel}{\level-1}%
    \foreach \i in {0,1,2} {
      \foreach \j in {0,1,2} {
        \pgfmathtruncatemacro{\skip}{(\i==1 && \j==1) ? 1 : 0}%
        \ifnum\skip=0
          \pgfmathsetmacro{\newx}{#1 + \i*\newsize}%
          \pgfmathsetmacro{\newy}{#2 + \j*\newsize}%
          \carpet{\newx}{\newy}{\newsize}{\nextlevel}%
        \fi
      }
    }
  \fi
}

\begin{tikzpicture}
  \foreach \stage in {0,1,2,3,4} {
    \begin{scope}[xshift=\stage*4cm]
      \draw[gray, thin] (0,0) rectangle (3,3);
      \carpet{0}{0}{3}{\stage}
      \node[below] at (1.5,-0.3) {Этап \stage};
    \end{scope}
  }
\end{tikzpicture}
\end{document}
\end{lstlisting}

\subsubsection{Результат}

\begin{center}
\begin{tikzpicture}[scale=0.35]
  \newcommand{\carpet}[4]{%
    \pgfmathtruncatemacro{\level}{#4}%
    \ifnum\level=0
      \fill[black] (#1,#2) rectangle +(#3,#3);
    \else
      \pgfmathsetmacro{\newsize}{#3/3}%
      \pgfmathtruncatemacro{\nextlevel}{\level-1}%
      \foreach \i in {0,1,2} {
        \foreach \j in {0,1,2} {
          \pgfmathtruncatemacro{\skip}{(\i==1 && \j==1) ? 1 : 0}%
          \ifnum\skip=0
            \pgfmathsetmacro{\newx}{#1 + \i*\newsize}%
            \pgfmathsetmacro{\newy}{#2 + \j*\newsize}%
            \carpet{\newx}{\newy}{\newsize}{\nextlevel}%
          \fi
        }
      }
    \fi
  }
  \foreach \stage in {0,1,2,3,4} {
    \begin{scope}[xshift=\stage*4cm]
      \draw[gray, thin] (0,0) rectangle (3,3);
      \carpet{0}{0}{3}{\stage}
      \node[below] at (1.5,-0.3) {\small Этап \stage};
    \end{scope}
  }
\end{tikzpicture}
\end{center}

\subsubsection{Анализ сложности}

На каждом уровне рекурсии количество квадратов увеличивается в 8 раз:
\begin{itemize}
  \item Этап 0: $8^0 = 1$ квадрат
  \item Этап 1: $8^1 = 8$ квадратов
  \item Этап 2: $8^2 = 64$ квадрата
  \item Этап 3: $8^3 = 512$ квадратов
  \item Этап 4: $8^4 = 4096$ квадратов
\end{itemize}

Временная сложность: $O(8^n)$, где $n$ — уровень рекурсии.

\section{Выводы}

В ходе выполнения работы были освоены следующие навыки:

\begin{enumerate}
  \item \textbf{Работа с узлами и стилями} — создание пользовательских стилей для узлов графа, использование полярных координат для размещения элементов.
  
  \item \textbf{Построение различных типов линий} — пунктирные, сплошные, кривые с параметрами кривизны, закругленные углы.
  
  \item \textbf{Визуализация математических функций} — использование команды \texttt{plot} для построения графиков функций, заданных аналитически.
  
  \item \textbf{Рекурсивное программирование в TikZ} — создание макросов с рекурсивными вызовами, условными операторами и циклами.
  
  \item \textbf{Работа с фракталами} — понимание принципов построения самоподобных структур, оптимизация рекурсивных алгоритмов.
\end{enumerate}

Пакет TikZ продемонстрировал высокую гибкость и мощность для создания сложных графических объектов. Особенно ценна возможность программного создания повторяющихся структур и интеграция с математическими вычислениями.

\section{Список использованных источников}

\begin{enumerate}
  \item Till Tantau. \textit{The TikZ and PGF Packages. Manual for version 3.1.10}. 2023.
  \item Benoit Mandelbrot. \textit{The Fractal Geometry of Nature}. W. H. Freeman and Company, 1982.
  \item Документация пакета TikZ: \url{https://ctan.org/pkg/pgf}
\end{enumerate}

\end{document}
